\section{A positive codimension-two bridge through qubit geometry}
Let $V\in\mathbb C^{n\times(n-2)}$ have unit-modulus entries and
$V^*V=nI_{n-2}$. Its residual Gram matrix
\[
 K=nI_n-VV^*
\]
is positive semidefinite of rank two, with diagonal entries two.
It has a factorization $K=RR^*$ with $R\in\mathbb C^{n\times2}$,
$R^*R=nI_2$, and every row of $R$ has squared norm two.
Define the real matrix
\[
 B_{ij}=\frac{|K_{ij}|^2}{2}-1.
\]

\begin{theorem}[An exact rank test for a two-column flat completion]
The matrix $B$ is a Gram matrix of unit vectors in $\mathbb R^3$.
The partial complex Hadamard matrix $V$ can be completed by two
unit-modulus columns if and only if $\operatorname{rank}B\le2$.
\end{theorem}
\begin{proof}
For each row $r_i$ of $R$, let $u_i=r_i^T/\sqrt2$ and form its unit
Bloch vector $b_i\in\mathbb R^3$. The elementary two-dimensional
identity $b_i\cdot b_j=2|u_i^*u_j|^2-1$ gives $B_{ij}$.
All other residual factors are obtained by a common two-dimensional
unitary change of basis. Such a factor has both entry moduli one
in every row exactly when every corresponding qubit has equal
coordinate probabilities, that is, every $b_i$ is perpendicular to
a fixed unit measurement axis. Such an axis exists exactly when
the Bloch vectors span dimension at most two. Explicitly, for a
unit normal $s$, diagonalize the traceless Hermitian matrix
$s\cdot\sigma$, whose eigenvalues are $\pm1$. Zero expectation in
each $u_i$ forces its two eigenbasis probabilities to be $1/2$.
This constructs the required residual factor.
\end{proof}

\paragraph{Formalization of the complex rank criterion.}
The independent \texttt{BlochChallenge} has two theorem targets and no
definition holes. The formalization constructs the actual residual
factor by complex orthonormal extension, proves its unit Bloch Gram
identity, derives the common normal by rank-nullity, and diagonalizes
the Hermitian axis to construct the two flat columns. The final
appended matrix retains every input entry and has full column Gram
$nI_n$. The converse is proved from the planar Bloch vectors of the
completing rows. The subsequent quaternary theorems are formalized in the independent
\texttt{QuaternaryChallenge}, with four theorem targets and no
definition holes. The automatic completion proof uses Gaussian parity
and an explicit two-axis Gram factorization.

\begin{theorem}[Quaternary input needs no extra phase freedom]
Suppose $n$ is even and $V$ has entries in $\mu_4$.
Then $B_{ij}\in\{-1,0,1\}$, and $\operatorname{rank}B\le3$.
Whenever $\operatorname{rank}B\le2$, the two completing columns
can also be chosen in $\mu_4$. Thus complex and quaternary
completability are equivalent for this codimension-two input.
\end{theorem}
\begin{proof}
The off-diagonal $K_{ij}$ are Gaussian integers with real and
imaginary parts of even sum, since they are sums of the even number
$n-2$ of fourth roots. Positive semidefiniteness gives $|K_{ij}|\le2$.
Thus their squared magnitudes are $0,2,4$, proving the assertion on
$B$. Distinct Bloch vectors are consequently parallel, antiparallel,
or orthogonal, and lie on at most three mutually orthogonal axes.

If the rank is at most two, rotate their plane to the equator.
After a common phase choice the residual rows have the form
$a_i(1,\delta_i)$ with $|a_i|=1$ and
$\delta_i\in\{1,-1,i,-i\}$. For two nonorthogonal rows,
\[
 K_{ij}=a_i\overline{a_j}(1+\delta_i\overline{\delta_j}).
\]
If the two $\delta$ values agree, the Gaussian integer $K_{ij}$
has modulus two, so division by two is a fourth root. If the axes
differ, its modulus is $\sqrt2$, so division by $1+i$ or $1-i$
is again a fourth root. Therefore $a_i\overline{a_j}\in\mu_4$
on every nonorthogonality edge. For rank two this graph is connected:
states on one axis are nonorthogonal to every state on the other.
A global phase makes all $a_i$ fourth roots. In rank one the two
antipodal groups are disconnected but mutually orthogonal; choose
their common phases independently without changing $K$.
This also makes every $a_i$ a fourth root. The resulting residual
matrix has entries in $\mu_4$ and the same Gram $K$, and appending
it to $V$ gives the required quaternary matrix.
\end{proof}
The proof is an exact algorithm using only $K$: take a reference
Bloch axis from a row of $B$, choose an orthogonal second axis if
needed, and label their antipodes by $1,-1,i,-i$. Along nonzero
$K_{ij}$ edges solve the displayed phase ratio, starting one phase
at one in each connected component. All recovered phases are in
$\mu_4$. This avoids numerical extraction of a Bloch rotation.

\begin{theorem}[Automatic quaternary completion]
Every even-order quaternary partial Hadamard matrix missing two
columns has a quaternary completion.
\end{theorem}
\begin{proof}
We show directly that three pairwise orthogonal Bloch vectors cannot
occur. For fourth roots $z_0,z_1,z_2$, the real part of
\[
 z_0\overline{z_1}+z_1\overline{z_2}+z_2\overline{z_0}
\]
is an odd integer. To see this, write $z_j=a_j+ib_j$ with
$a_j+b_j$ odd. Modulo two its real part is
$\sum_{\mathrm{cyc}}[a_ja_{j+1}+(1-a_j)(1-a_{j+1})]=1$.
Summing this identity over the even number $n-2$ of given columns
shows that, for distinct row indices $i,j,k$,
\[
 \operatorname{Re}(K_{ij}+K_{jk}+K_{ki})\equiv0\pmod2.
\]
If the corresponding three Bloch vectors were pairwise orthogonal,
then $B_{ij}=B_{jk}=B_{ki}=0$. Consequently each of the Gaussian
integers $K_{ij},K_{jk},K_{ki}$ would have squared modulus two,
and hence odd real part. Their sum would have odd real part,
a contradiction.

Unit vectors with pairwise inner products in $\{-1,0,1\}$ are
parallel, antiparallel, or orthogonal. Choose one reference vector.
If every vector is parallel to it, their span has dimension one.
Otherwise choose a perpendicular second vector. Any vector not
parallel to either would be perpendicular to both, producing the
forbidden triple. Thus the span has dimension at most two, so
$\operatorname{rank}B\le2$. The preceding fourth-root completion
criterion supplies the two completing columns.
\end{proof}

This proof uses Gaussian parity and the preceding Bloch criterion;
it does not require a real codimension-four completion theorem.
Classical completion results include E. Verheiden,
\emph{Integral and rational completions of combinatorial matrices},
J. Combinatorial Theory A 25 (1978), 267--276, and his dissertation
\url{https://thesis.caltech.edu/2899/}. The latter records Hall's
four-row theorem and its extension to seven.
The codimension-one statement is simpler: its residual has rank one
and diagonal one, and one of its columns directly supplies the
missing fourth-root column up to phase. Neither result supplies
the missing large partial frame at an unknown order.

The checker verifies the Bloch rank, phase-graph algorithm, and exact
Gaussian Gram identities at orders four, six and eight. A bounded
search of small quaternary partial frames found no rank-three input,
consistent with the classical-completion consequence. A three-axis
abstract Gaussian residual is also tested to distinguish the
residual geometry from a realizable quaternary partial frame.
